% !TeX root = ../../Book3.tex
% 纲要标题：### G.2 实根理想与实零点定理

\section{实根理想与实零点定理}
\label{b3:sec:G02}

在实闭域上，平方和为零迫使每个平方项为零，这给普通根理想增加了一类关系。实根理想与实零点的消失理想之间的对应，正是把这些额外关系代数化的结果。

% 纲要细目（与计划纲要同步）：
% * 比较原理想、普通根理想、实根理想与实消失理想；平方和方程正式说明点集相同而理想不同
% * 由不可行性证书、局部化及偶次幂清分母完整证明实零点定理；本节向前引用 G.4 而不构成循环
% * 实谱的点为素理想及其剩余分式域上的序；同一素理想可给不同实谱点
% #### G.2.1 实理想、实根理想与实零点集
% #### G.2.2 平方和方程的实消失理想
% #### G.2.3 实零点定理与实谱的基本对象

\subsection{实理想、实根理想与实零点集}
\label{b3:subsec:G02:01}

若只取普通根理想，实零点仍可能丢失信息。例如 \((x^2+y^2)\subseteq\mathbb R[x,y]\) 是根理想，但它的实零点只有原点，消失理想为 \((x,y)\)。缺少的信息来自“平方和为零迫使每项为零”。

\begin{definition}\label{b3:def:real-ideal}
设 \(A\) 为交换环，\(I\subseteq A\) 为理想。若对任意正整数 \(r\) 和任意 \(a_1,\ldots,a_r\in A\)，只要 \(a_1^2+\cdots+a_r^2\in I\) 就有 \(a_1,\ldots,a_r\in I\)，则称 \(I\) 为\term{实理想}[real ideal]。所有包含 \(I\) 的实理想（包括 \(A\) 本身）的交称为 \(I\) 的\term{实根理想}[real radical]，记为 \(\sqrt[\mathrm r]{I}\)。
\end{definition}
\symidx{\(\sqrt[\mathrm r]{I}\)：理想 \(I\) 的实根理想}

实理想必为根理想：若 \(a^N\in I\)，取 \(2^e\ge N\)，则 \(a^{2^e}\in I\)，反复用单个平方的条件得 \(a\in I\)。实理想的交仍为实理想，因此上一定义确实给出包含 \(I\) 的最小实理想。

\begin{proposition}\label{b3:prop:real-vanishing-ideal}
设 \(R\) 为实闭域，\(S\subseteq R^n\)，并记 \(I(S)=\{f\in R[x_1,\ldots,x_n]:f(a)=0\;\text{对所有}\;a\in S\}\)。则消失理想 \(I(S)\) 是实理想。
\end{proposition}
\begin{proof}
若 \(\sum_j f_j^2\) 在 \(S\) 上为零，则对每个 \(a\in S\)，非负数之和 \(\sum_jf_j(a)^2=0\) 迫使各 \(f_j(a)=0\)。故各 \(f_j\in I(S)\)。空集情形给出整个多项式环，也满足结论。
\end{proof}

对实闭域 \(R\) 上的多项式理想 \(I\subseteq R[x_1,\ldots,x_n]\)，普通根理想、实根理想与实零点的消失理想之间的关系是
\[
I\subseteq\sqrt I\subseteq\sqrt[\mathrm r]{I}=I(V_R(I)).
\]
第一个包含由普通根理想的定义得到，第二个来自实理想必为根理想；上面的命题还给出 \(\sqrt[\mathrm r]{I}\subseteq I(V_R(I))\)。最后的等号由\cref{b3:thm:real-nullstellensatz}证明，它把实点上的消失条件转成有限平方和证书。

\subsection{平方和方程的实消失理想}
\label{b3:subsec:G02:02}

\begin{proposition}\label{b3:prop:sum-square-real-zero}
设 \(R\) 为实闭域，\(n,s\ge1\)，\(f_1,\ldots,f_s\in R[x_1,\ldots,x_n]\)。记 \(V_R(f_1,\ldots,f_s)\) 为这些多项式在 \(R^n\) 中的共同零点集，\(I(S)\) 为在集合 \(S\subseteq R^n\) 上处处为零的多项式组成的理想。则
\[
V_R(f_1^2+\cdots+f_s^2)=V_R(f_1,\ldots,f_s).
\]
特别地，\(I(V_R(x_1^2+\cdots+x_n^2))=(x_1,\ldots,x_n)\)。
\end{proposition}
\begin{proof}
第一式仍由各平方非负得到。第二式中零点只有原点；一个多项式在原点为零当且仅当其常数项为零，这又等价于属于 \((x_1,\ldots,x_n)\)。
\end{proof}

具体地，取 \(I=(x^2+y^2)\subseteq R[x,y]\)，则
\[
\sqrt I=(x^2+y^2)\subsetneq(x,y)=\sqrt[\mathrm r]{I}.
\]
\(x^2+y^2\) 在实闭域 \(R\) 上不可约，而 \(R[x,y]\) 是唯一分解整环，所以 \(I\) 为素理想，已经等于普通根理想。另一方面，每个含 \(I\) 的实理想都含 \(x,y\)；\((x,y)\) 又是原点的实消失理想，故它恰为 \(\sqrt[\mathrm r]{I}\)。这直接显示第二个包含可以严格。

\ProseParagraphBreak
这里把普通理想换成实根理想，不是在环中增加一个等式 \(x_1=0\)，而是在指定实闭域的点上辨认哪些等式已经由平方和关系强制成立。相反，复数域上 \(x^2+y^2=(x+\mathrm iy)(x-\mathrm iy)\) 的零点包含两条直线。基础域的选择改变了几何对象。

\ProseParagraphBreak
另一个极端是 \(I=(x^2+1)\subseteq R[x]\)。它没有实零点，且 \(1^2+x^2\in I\)。任何包含它的实理想都含 \(1\)，所以 \(\sqrt[\mathrm r]{I}=R[x]\)。这比仅知道 \(I\) 为极大理想更贴近实零点问题。

\subsection{实零点定理与实谱的基本对象}
\label{b3:subsec:G02:03}

\begin{theorem}[实零点定理]\label{b3:thm:real-nullstellensatz}
设 \(R\) 为实闭域，\(A=R[x_1,\ldots,x_n]\)，\(I\subseteq A\) 为理想。记 \(V_R(I)\) 为 \(I\) 在 \(R^n\) 中的共同零点集，\(I(V_R(I))\) 为在该集合上处处为零的多项式组成的理想。则
\[
I(V_R(I))=\sqrt[\mathrm r]{I}
=\{f\in A:\ f^{2m}+h_1^2+\cdots+h_s^2\in I
\text{，某}\;m\ge1,\ h_j\in A\}.
\]
\end{theorem}

本定理的证明使用\subsecref{b3:subsec:G04:02}中证明的 Krivine--Stengle 型不可行性证书\cref{b3:thm:real-infeasibility-certificate}，这里取没有不等式约束的情形，其预序就是有限平方和集合。该证书的证明由极大预序、实闭包与本附录已证的量词消去得到；读完那个证明后，即可完成下面实零点定理的推导。

\ProseParagraphBreak
具体方法与\subsecref{b3:subsec:0603:02}的 Rabinowitsch 技巧相同：加入 \(Zf=1\)，把 \(f\) 在原零点集上恒为零转成增广系统无实解；不可行性证书给出 \(-1\) 的平方和及理想成员表示；再在局部化中令 \(Z=f^{-1}\)，用偶次幂清去分母，得到 \(f^{2m}\) 加平方和属于 \(I\) 的恒等式。与代数闭域的版本相比，这里无解性给出的是平方和证书，清分母时必须保存这些平方项。

\begin{proof}
若 \(f=0\)，取 \(m=1\)、平方和为零便有证书。以下设 \(f\ne0\) 且 \(f\in I(V_R(I))\)。由\cref{b3:cor:hilbert-module-consequences}，域 \(R\) 上的有限变量多项式环 \(A\) 是 Noether 环，可取 \(I\) 的有限组生成元。增加一个变量 \(Z\) 后，这些生成元等于零并且 \(1-Zf=0\) 的有限系统没有实解。将\cref{b3:thm:real-infeasibility-certificate}用于这个系统，得到
\[
-1=\sum_j h_j(x,Z)^2+u(x,Z),
\qquad u\in I A[Z]+(1-Zf).
\]
把 \(u\) 写成有限和
\[
u=\sum_{\nu=1}^t c_\nu(x,Z)a_\nu(x)+(1-Zf)v(x,Z),
\qquad a_\nu\in I.
\]
在局部化 \(A_f\) 中代入 \(Z=f^{-1}\)，得到
\[
-1-\sum_j h_j(x,f^{-1})^2
=\sum_\nu c_\nu(x,f^{-1})a_\nu\in IA_f.
\]
这不是在 \(f=0\) 的点求值。取 \(m\ge1\) 足够大，使每个 \(f^m h_j(x,f^{-1})\) 与 \(f^{2m}c_\nu(x,f^{-1})\) 都属于 \(A\)；例如同时要求 \(m\ge\deg_Zh_j\)、\(2m\ge\deg_Zc_\nu\) 即可。这样平方项与理想成员的系数都用同一个偶次幂清去分母，得到
\[
-f^{2m}-\sum_j\bigl(f^m h_j(x,f^{-1})\bigr)^2
=\sum_\nu\bigl(f^{2m}c_\nu(x,f^{-1})\bigr)a_\nu\in I.
\]
最后利用 \(I=-I\)，把整个等式乘以 \(-1\)，才得到所需的正号形式
\[
f^{2m}+\sum_j\bigl(f^m\cdot h_j(x,f^{-1})\bigr)^2\in I.
\]
各括号此时都属于 \(A\)。

记全部具有上述证书的元素所成集合为 \(C\)。在任意含 \(I\) 的实理想 \(J\) 中，证书迫使 \(f^m\in J\)，而实理想为根理想，故 \(f\in J\)。所以 \(C\subseteq\sqrt[\mathrm r]{I}\)。另一方面，\cref{b3:prop:real-vanishing-ideal}说明 \(I(V_R(I))\) 本身是含 \(I\) 的实理想，因而
\(\sqrt[\mathrm r]{I}\subseteq I(V_R(I))\)。第一段已经证明 \(I(V_R(I))\subseteq C\)，三个包含合起来即得等式。
\end{proof}

实理想版本见\cite[Appendix A, Theorem A.5.15]{Mangolte2020RealAlgebraicVarieties}；\(\mathbb R\) 上的证书版本见\cite[Section 7.4, Theorem 7.5]{Sturmfels2002PolynomialEquations}。证书是有限多项式恒等式，能够用精确运算复查。

\begin{table}[htbp]
\centering
\caption{普通根理想与实根理想}
\label{b3:tab:G-radical-comparison}
% !TeX root = ../Book3.tex
\begin{tabularx}{\textwidth}{@{}>{\raggedright\arraybackslash}p{.20\textwidth}>{\raggedright\arraybackslash}p{.33\textwidth}>{\raggedright\arraybackslash}X@{}}
\toprule
对象 & 代数条件或定义 & 零点解释 \\
\midrule
普通根理想 \(\sqrt I\) & 某个 \(m\ge1\) 使 \(f^m\in I\) & 在代数闭域上，它等于共同零点集的消失理想 \\
实根理想 \(\sqrt[\mathrm r]{I}\) & 某个 \(m\ge1\) 与有限列 \(h_j\) 使 \(f^{2m}+\sum_jh_j^2\in I\) & 在实闭域上，它等于共同实零点集的消失理想 \\
实消失理想 \(I(V_R(I))\) & \(f(a)=0\) 对每个 \(a\in V_R(I)\) 成立 & 定义来自实点；实零点定理将其转为上一行的有限证书 \\
\bottomrule
\end{tabularx}

\end{table}

\cref{b3:tab:G-radical-comparison}分别在代数闭域和实闭域上比较零点的消失理想。代数闭域的零点由普通根理想描述；实闭域的零点还会强制平方和中的每个平方项为零，因此需要实根理想。

\ProseParagraphBreak
若要同时考察不同实闭扩张中的点，就会用到\term{实谱}[real spectrum]。对交换环 \(A\)，其点可定义为一对 \((\mathfrak p,\le)\)：\(\mathfrak p\) 是素理想，\(\le\) 是域 \(\operatorname{Frac}(A/\mathfrak p)\) 上的序。由条件 \(f_1>0,\ldots,f_r>0\) 给出基本开集。取实点 \(a\) 得到 \(\mathfrak p=\ker(A\to R)\)，并继承相应的序；但一般实谱点不必来自 \(R^n\) 中的某个点。这里的附加数据是序，不能把实谱直接等同于普通素谱的某个无附加结构子集。

\ProseParagraphBreak
例如 \(A=\mathbb R[x]\) 的同一个素理想 \(\mathfrak p=(0)\) 可以配上不同的序。对其剩余分式域 \(\mathbb R(x)\)，记 \(\operatorname{lc}(u)\) 为非零多项式 \(u\) 的最高次项系数，规定
\[
\frac{u(x)}{v(x)}>_+0
\quad\Longleftrightarrow\quad
\operatorname{lc}(u)\operatorname{lc}(v)>0.
\]
交叉相乘说明这与分式的表示无关。正分式可写成 \(uv/v^2\)，其分子最高次项系数为正；这种表示下，两个正分式的和与积仍有正的最高次项系数，因而上述规则确实定义域序。在这个序中 \(x-c>_+0\) 对每个 \(c\in\mathbb R\) 成立。再用自同构 \(x\mapsto-x\) 拉回该序，得到另一个序 \(\le_-\)，其中 \(x<c\) 对每个实常数 \(c\) 成立。两者都延伸通常的实数序，却给 \(x\) 相反的符号：\((\mathfrak p,\le_+)\) 与 \((\mathfrak p,\le_-)\) 是两个不同的实谱点。素理想决定代数支集，序还记录这份支集上元素的符号。
